\chapter{Internal Game Structure: Logic Behind Counter-Strike 2}

\section{Understanding Entity Lists in Video Games}

In modern game engines—particularly those designed for real-time, multiplayer environments—an \textbf{entity list} is a fundamental data structure responsible for managing all interactive objects within the game world. These entities include players, non-player characters (NPCs), projectiles, grenades, environmental props, and more.

\begin{tcolorbox}[colback=gray!10, colframe=gray!50, sharp corners, boxrule=0pt]
\textbf{What is an Entity?} \\
An \textit{entity} is any dynamic object in the game world, such as a player, NPC, projectile, or interactive prop, managed and updated by the engine each frame.
\end{tcolorbox}

At the core, each entity is an instance of a class containing essential data such as position, health, team ID, state flags, and animation or input components. The engine maintains and updates these entities each frame to simulate gameplay logic, render visuals, process collisions, and evaluate state-based behavior such as death, visibility, or interaction.

\subsection{Common Structures for Entity Lists}

Entity lists are implemented using various structures, depending on the performance needs and memory model of the engine:

\begin{itemize}
    \item \textbf{\texttt{Array of Entity Objects}}: \\
    A flat, contiguous layout offering fast access, but limited flexibility.
    
    \item \textbf{\texttt{Array of Pointers}}: \\
    A static array referencing dynamically allocated entities. Frequently used in older engine architectures.
    
    \item \textbf{\texttt{std::vector}}: \\
    A dynamic container ideal for runtime creation and destruction of entities.
    
    \item \textbf{\texttt{Linked List}}: \\
    A flexible memory structure where each node points to the next entity, at the cost of linear traversal time.
    
    \item \textbf{\texttt{HashMap / Tree}}: \\
    Enables fast lookup and categorization of entities based on IDs or class types.
\end{itemize}

\begin{tcolorbox}[colback=gray!10, colframe=gray!50, sharp corners, boxrule=0pt]
    \textbf{What is the Paged Entity System?} \\
    A memory management technique where entities are grouped into fixed-size pages, improving scalability, memory safety, and preventing invalid accesses when entities are destroyed or recreated.
\end{tcolorbox}
    
While earlier iterations of the Source engine (such as CS:GO) utilized a simple flat array of entity pointers for entity management, the modern Source 2 engine in Counter-Strike 2 introduces a \textbf{paged entity system}. This design organizes entities into memory pages and utilizes safe handle-based referencing, offering greater scalability, safety, and efficiency across dynamically allocated objects.    

\section{Entity Organization in Counter-Strike 2}

In \textit{Counter-Strike 2}, every interactive object—whether a player, weapon, projectile, or environmental element—is internally represented as a dynamic \textbf{entity}. These entities are managed using unique \textit{handles}, lightweight references that combine an entity index and a generation counter to ensure safe access and lifecycle management.

\begin{tcolorbox}[colback=gray!10, colframe=gray!50, sharp corners, boxrule=0pt]
\textbf{What is a Handle?} \\
A \textit{handle} is a compact reference to an entity, typically composed of an index and a generation counter. It allows efficient and safe resolution of entities, preventing issues such as dangling pointers when entities are destroyed and recreated.
\end{tcolorbox}

This handle-based system enables robust referencing, safeguarding against invalid memory access even in highly dynamic environments where entities frequently appear and disappear during gameplay.

To map a handle to an actual entity in memory, the Source 2 engine employs a \textbf{two-level paging system}—a mechanism conceptually similar to virtual memory in modern operating systems. This approach balances memory efficiency, scalability, and safety, particularly as entity counts grow in large matches or complex game scenarios.

When analyzing player-specific data structures, two core classes emerge as critical pillars:

\begin{itemize}
    \item \textbf{CBasePlayerController} — Manages high-level metadata about a player, including their sanitized name, Steam ID, current input state, and a handle linking to their in-world representation (the pawn).
    \item \textbf{C\_CSPlayerPawn} — Represents the player's physical presence within the game world, storing vital gameplay attributes such as position, health, team affiliation, view angles, and animation states.
\end{itemize}

Both player-related classes inherit from \textbf{C\_BaseEntity}, the core base class that provides common properties, behaviors, and networking support for all entities in the engine.

The following are simplified and annotated representations of these structures:

\vspace*{0.5cm}

\noindent\textbf{1. Base Entity – \texttt{C\_BaseEntity}}

\begin{lstlisting}[style=cppstyle, caption={C\_BaseEntity – Core Class for All Entities}]
/*********************************************
 * C_BaseEntity:
 * Base class for all entities in Source 2.
 *********************************************/
class C_BaseEntity {
public:
    CHandle<C_BaseEntity>     m_hEffectEntity;    // Effect entity 
    CHandle<C_BaseEntity>     m_hOwnerEntity;     // Owner 
    CHandle<C_BaseEntity>     m_hGroundEntity;    // Entity standing on 
    CBodyComponent::Storage_t m_CBodyComponent;   // Physics
    int32_t                   m_iHealth;          // Entity health
    Vector3                   m_vPosition;        // World-space position

    // ... Additional fields (e.g., team, animation state)
};
\end{lstlisting}


\noindent\textbf{2. Player Controller – \texttt{CBasePlayerController}}

\begin{lstlisting}[style=cppstyle, caption={CBasePlayerController – Player Metadata and Input Handler}]
/********************************************************
 * CBasePlayerController:
 * Manages metadata, Steam ID, input, and pawn handle.
 ********************************************************/
class CBasePlayerController : public C_BaseEntity {
public:
    std::string              m_sSanitizedPlayerName;  // Player name 
    uint64_t                 m_steamID;               // Steam Id
    InputState_t             m_inputState;            // Current input 
    CHandle<C_CSPlayerPawn>  m_hPawn;                  // Handle to pawn 
    uint32_t                 m_iTeamNum;              // Team ID
    bool                     m_bIsConnected;          // Connection status

    // ... Other metadata (e.g., score, latency)
};
\end{lstlisting}

\newpage


\noindent\textbf{3. Player Pawn – \texttt{C\_CSPlayerPawn}}

\begin{lstlisting}[style=cppstyle, caption={C\_CSPlayerPawn – In-World Representation of the Player}]
/********************************************************
 * C_CSPlayerPawn:
 * Represents the physical player in the world.
 ********************************************************/
class C_CSPlayerPawn : public C_BaseEntity {
public:
    Vector3            m_vPosition;      // World position
    uint32_t           m_iHealth;         // Health
    uint32_t           m_iTeamNum;        // Team ID
    QAngle             m_angViewAngles;   // View angles
    AnimationState_t   m_AnimState;       // Animation state

    // ... Other attributes
};
\end{lstlisting}


This class hierarchy enables a clear separation of concerns between the \textit{logical control} of the player (controller) and their \textit{physical simulation} within the game world (pawn). This design also accommodates edge cases such as spectator modes, bot logic, and replay systems, where a controller may exist without an active pawn.

Moreover, the architecture follows a variant of the \textbf{Entity-Component-System (ECS) paradigm}, promoting modularity and flexibility. This pattern is especially beneficial for external tools such as mods, analytics systems, and cheat frameworks.

\section{Entity List Traversal}

Having established the structural foundations of player entities, the next step involves accessing and enumerating these entities during runtime. This process—known as \textbf{entity list traversal}—is essential for real-time game modification, data extraction, and the development of external tools.

\subsection{Locating and Identifying the Entity List in Memory}

In the absence of \textbf{debugging symbols} or \textbf{SDK access}—a common scenario in protected or closed-source games—discovering the location and structure of the \textbf{entity list} becomes a crucial part of the reverse engineering process. Two primary approaches are typically employed: \textit{\textbf{dynamic memory inspection}} and \textit{\textbf{static binary analysis}}.


These complementary methods allow external programs to interact systematically with the game’s internal structures.


\begin{tcolorbox}[colback=gray!03, colframe=black!15, sharp corners, boxrule=0.3mm, width=\textwidth]

    \subsubsection*{Dynamic Approach: Real-Time Memory Inspection}
    
    Dynamic memory inspection focuses on observing changes in the game’s memory while it is running, providing a practical way to locate the entity list without dissecting the entire executable.
    
    The discovery process generally follows these steps:
    
    \begin{enumerate}
        \item \textbf{Scan for the Local Player’s Position:}  
        Search for changing float values that represent the player's coordinates (\texttt{X}, \texttt{Y}, \texttt{Z}) during movement.
    
        \item \textbf{Filter Out Non-Relevant Addresses:}  
        Discard addresses that stop updating when the game is paused or unfocused, as they typically reflect cached or rendering-only data.
    
        \item \textbf{Analyze Accessing Instructions:}  
        Use memory inspection tools to detect which instructions interact with these addresses.
    
        \item \textbf{Detect Multi-Entity Patterns:}  
        Identify if the same instruction accesses positions of multiple players, suggesting iteration over an entity container.
    
        \item \textbf{Recover Entity Base Pointers:}  
        Subtract known field offsets (e.g., \texttt{m\_vPosition}) from accessed addresses to obtain base pointers to entities.
    
        \item \textbf{Locate the Entity List Structure:}  
        Search for contiguous memory regions holding these base pointers, checking for uniform spacing, characteristic of structured storage (arrays or paged layouts).
    
        \item \textbf{Confirm and Generalize:}  
        Build a traversal loop using pointer arithmetic to iterate through the list and access relevant entity attributes such as health, team ID, and state flags.
    \end{enumerate}
    
    
    \begin{tcolorbox}[colback=gray!07, colframe=black!30, sharp corners, boxrule=0pt]
    \textbf{Important:}  
    Not all structures resembling entity lists are gameplay-relevant. Careful validation is necessary to avoid mistaking render-only data for true gameplay entities.
    \end{tcolorbox}
    
\end{tcolorbox}
    
    

\subsubsection*{Types of Entity Lists Encountered}

During memory analysis, it is typical to encounter two distinct categories of entity lists:

\begin{itemize}
    \item \textbf{Render-Only Lists:}  
    Contain data purely used for graphical interpolation or rendering, but lack gameplay-critical attributes.

    \item \textbf{Gameplay-Relevant Lists:}  
    Include full gameplay state data such as health, team identifiers, ammunition counts, and behavioral flags.
\end{itemize}

Locating a render-only list is relatively quick, but isolating a gameplay-relevant entity list demands deeper inspection and runtime validation.

\begin{tcolorbox}[colback=gray!03, colframe=black!15, sharp corners, boxrule=0.3mm, width=\textwidth]

    \subsubsection*{Static Approach: Reverse Engineering the Binary}
    
    When dynamic analysis alone is insufficient or incomplete, static binary analysis becomes essential. This method involves disassembling and inspecting the game executable to directly uncover the internal logic responsible for managing entities.
    
    \vspace{0.4cm}
    
    \textbf{Common tools used for static analysis:}
    
    \begin{itemize}
        \item \textbf{IDA Pro}
        \item \textbf{Ghidra}
        \item \textbf{Binary Ninja}
    \end{itemize}
    
    \vspace{0.2cm}
    
    Static analysis provides critical insights such as:
    
    \begin{itemize}
        \item Identification of virtual tables (\textbf{vtables}) associated with entity classes.
        \item Extraction of Run-Time Type Information (\textbf{RTTI}) metadata, revealing class hierarchies and inheritance relationships.
        \item Detection of key instruction patterns that govern entity list traversal and access.
    \end{itemize}
    
    \vspace{0.2cm}
    
    \textbf{IDA Pro Plugins that enhanced the process:}
    
    \begin{itemize}
        \item \textbf{ClassInformer}:  
        Automatically locates and reconstructs RTTI structures, making it easier to understand class inheritance and find classes like \texttt{C\_BaseEntity} and its derivatives.
    
        \item \textbf{StringAssociator}:  
        Helps associate readable strings (such as entity names or debug labels) with their corresponding functions or objects, facilitating navigation through the binary.
    
        \item \textbf{SigMaker}:  
        Creates byte-pattern signatures for functions or structures. These signatures are essential to reliably relocate functionality across game updates where absolute addresses change.
    \end{itemize}
    
    \vspace{0.4cm}
    
    This hybrid strategy—leveraging both real-time memory inspection and deep static disassembly—has been widely adopted and discussed in the reverse engineering community.  
    An excellent case study can be found in a collaborative thread on the UnknownCheats forum\footnote{\url{https://www.unknowncheats.me/forum/general-programming-and-reversing/564058-entity-list.html}}.
    
    \vspace{0.4cm}
    
    \noindent\textbf{In the next section, we will explain how we extracted the internal logic behind entity list resolution through careful static analysis of the game’s binary.}
    
\end{tcolorbox}
    
\subsection{Reverse Engineering the Engine's Entity Resolution Logic}

By analyzing the \texttt{client.dll} binary of CS2 using static reverse engineering techniques, we can uncover the specific function responsible for resolving entity handles into entity pointers.

The following pseudocode extracted from Ghidra outlines the resolution logic:

\begin{lstlisting}[style=cppstyle, caption={Accessing an Entity in the Entity List}]
Entity_t *EntityList(longlong EntityListPointer, uint EntityListIndex)
{
    Entity_t *Entity;
    longlong EntitySublist;

    if ((((EntityListIndex < 32767) && (EntityListIndex >> 9 < 64)) &&
         (EntitySublist = *(EntityListPointer + 16 + (EntityListIndex >> 9) * 8), EntitySublist != 0))
        && ((Entity = EntitySublist + (EntityListIndex & 511) * 120, Entity != NULL &&
             ((*(Entity + 2) & 32767) == EntityListIndex))))
    {
        return *Entity;
    }
    return 0;
}   
\end{lstlisting}

\noindent
The process of resolving an entity from its handle follows a clear, sequential logic:

\begin{enumerate}
    \item \textbf{Validate the Index:}  
    Ensure that the provided \texttt{EntityListIndex} is valid:
    \begin{itemize}
        \item \texttt{EntityListIndex < 32767} — The index must be within the allowable entity range.
        \item \texttt{EntityListIndex >> 9 < 64} — The index must map to one of 64 available pages.
    \end{itemize}

    \item \textbf{Find the Correct Page:}  
    Calculate the address of the memory page (sublist) containing the entity:
    \begin{itemize}
        \item \texttt{EntitySublist = *(EntityListPointer + 16 + (EntityListIndex >> 9) * 8)}
    \end{itemize}
    If the page is invalid (\texttt{EntitySublist == 0}), the search stops here.

    \item \textbf{Locate the Entity Within the Page:}  
    Find the exact entity inside the page:
    \begin{itemize}
        \item \texttt{Entity = EntitySublist + (EntityListIndex \& 511) * 120}
    \end{itemize}
    This isolates the local index inside the page (modulo 512) and multiplies by the entity size (120 bytes).

    \item \textbf{Verify the Entity:}  
    Before trusting the result, two conditions must be satisfied:
    \begin{itemize}
        \item \texttt{Entity != NULL} — The pointer must be valid.
        \item \texttt{(*(Entity + 2) \& 32767) == EntityListIndex} — The stored index must match the requested one.
    \end{itemize}

    \item \textbf{Return the Result:}  
    If all checks pass, return the pointer to the entity.  
    Otherwise, return \texttt{0}.
\end{enumerate}

\vspace{0.5cm}

\begin{tcolorbox}[colback=gray!05, colframe=gray!50, sharp corners, boxrule=0pt]
\textbf{Summary:}  
The handle resolution process ensures safety by validating the index, finding the correct memory page, isolating the entity within it, and double-checking its integrity before returning a valid pointer.
\end{tcolorbox}

\subsubsection{Efficient Entity Tracking: Hooking \texttt{OnAddEntity} and \texttt{OnRemoveEntity}}

While it is theoretically possible to retrieve all active entities by continuously iterating through the entire entity list each frame (or each game tick), this approach carries significant performance drawbacks.  
Given that the entity list can accommodate up to \textbf{32,766 entities} simultaneously, a full traversal on every tick would introduce substantial overhead—especially when injecting an external cheat into the game process.

During our research, we successfully \textbf{identified and reconstructed} the internal entity list traversal logic used by the engine. We have integrated this traversal into our cheat, allowing us to programmatically enumerate all active entities when needed.

However, rather than calling this loop every frame—which would severely impact performance—we apply the following strategy:

\begin{itemize}
    \item Upon joining a match, we execute our custom traversal \textbf{once} to perform an initial discovery and caching of all existing entities.
    \item After this one-time initialization, we \textbf{hook} two critical engine callbacks:
    \begin{itemize}
        \item \texttt{OnAddEntity} — Triggered whenever a new entity is created or spawned.
        \item \texttt{OnRemoveEntity} — Triggered whenever an entity is deleted or removed.
    \end{itemize}
\end{itemize}

This hybrid approach allows us to dynamically maintain an accurate and up-to-date internal list of game entities without imposing unnecessary overhead on each frame.
